\chapter{Marco teórico}\label{cap:teoria}

Este capítulo reúne la teoría que se usa en el dimensionamiento. Se parte del caso más
simple (equilibrio de fuerzas) y se agregan, en orden, el comportamiento del flujo a
través del disco, el de cada sección de pala, el del perfil a bajo número de Reynolds y
la dinámica de las palas articuladas.

%=====================================================================
\section{Equilibrio vertical}
A velocidad de descenso constante $V$, el peso iguala la suma de las fuerzas que frenan:
rotor, drogue y cuerpo.
\begin{equation}
  W = m g = \tfrac12 \rho V^2 \left( \CR A_R + C_{D,d} A_d + C_{D,c} A_c \right)
  \;\Rightarrow\;
  V = \sqrt{\dfrac{2 m g}{\rho \left( \CR A_R + C_{D,d} A_d + C_{D,c} A_c\right)}}
  \label{eq:equilibrio}
\end{equation}
donde $A_R=\pi R^2$ es el área del disco del rotor, $\CR$ su coeficiente de resistencia
equivalente, $A_d$ y $C_{D,d}$ los del drogue y $A_c$ y $C_{D,c}$ los del cuerpo. Todo el
problema se reduce a conocer $\CR$, lo que desarrollan las secciones siguientes.

%=====================================================================
\section{Teoría de cantidad de movimiento en vuelo axial}\label{sec:momento}

\subsection{Vuelo estacionario y velocidad inducida de referencia}
Un rotor que sostiene un empuje $T$ acelera el aire a través del disco. En vuelo
estacionario, la velocidad inducida en el disco es~\cite{leishman2006,johnson1980}
\begin{equation}
  \vh = \sqrt{\frac{T}{2\rho A_R}} ,
  \label{eq:vh}
\end{equation}
que se usa como escala para todos los regímenes axiales.

\subsection{Ramas válidas de la teoría}
Con velocidad axial $V_c$ (positiva subiendo, negativa bajando) e inducida $\vi$
(positiva hacia abajo), la teoría de cantidad de movimiento da dos ramas con solución
física:
\begin{align}
  \text{ascenso } (V_c \ge 0):&\quad \frac{\vi}{\vh} = -\frac{V_c}{2\vh} +
      \sqrt{\left(\frac{V_c}{2\vh}\right)^2 + 1},\label{eq:ascenso}\\
  \text{molino frenante } (V_c \le -2\vh):&\quad \frac{\vi}{\vh} = -\frac{V_c}{2\vh} -
      \sqrt{\left(\frac{V_c}{2\vh}\right)^2 - 1}.\label{eq:molino}
\end{align}
Entre $-2\vh < V_c < 0$ el flujo no tiene una dirección definida a través del disco: es
la zona del \textbf{anillo de vórtices} y de la \textbf{estela turbulenta}, donde la
teoría no vale y se usan datos experimentales~\cite{castles1951,johnson2005,leishman2004}.

\subsection{Curva empírica en descenso}
Leishman~\cite{leishman2006} ajusta los datos de descenso axial con un polinomio,
válido para $-2\le V_c/\vh\le0$:
\begin{equation}
  \frac{\vi}{\vh} = \kappa + k_1 x + k_2 x^2 + k_3 x^3 + k_4 x^4, \qquad x=\frac{V_c}{\vh},
  \label{eq:poly}
\end{equation}
con $\kappa=1{,}15$, $k_1=-1{,}125$, $k_2=-1{,}372$, $k_3=-1{,}718$, $k_4=-0{,}655$.%
\footnote{Coeficientes tomados de Leishman (2006), cap.~2; conviene verificarlos contra la
edición disponible. El resultado del capítulo~\ref{cap:dim} se contrasta además con un
modelo de elemento de pala independiente.}

\begin{figure}[H]
\centering
\begin{tikzpicture}
\begin{axis}[width=0.92\textwidth, height=8.2cm,
  xmin=-4, xmax=1.5, ymin=0, ymax=3.6,
  xlabel={Velocidad axial $V_c/\vh$ \quad(negativa: descenso)},
  ylabel={Velocidad inducida $\vi/\vh$},
  legend style={at={(0.02,0.98)}, anchor=north west, font=\footnotesize, cells={anchor=west}},
  clip=false]
  % zonas
  \fill[cfreno!8] (axis cs:-2,0) rectangle (axis cs:-1.4,3.6);
  \fill[ccuerda!10] (axis cs:-1.4,0) rectangle (axis cs:0,3.6);
  \fill[cfijo!8] (axis cs:-4,0) rectangle (axis cs:-2,3.6);
  % ramas de cantidad de movimiento
  \addplot[cfijo, very thick, domain=0:1.5, samples=60] {-x/2 + sqrt(x^2/4+1)};
  \addlegendentry{Teoría: ascenso}
  \addplot[cfijo, very thick, dashed, domain=-4:-2, samples=80] {-x/2 - sqrt(max(x^2/4-1,0))};
  \addlegendentry{Teoría: molino frenante}
  % polinomio empírico
  \addplot[crot, very thick, domain=-2:0, samples=80]
    {1.15 - 1.125*x - 1.372*x^2 - 1.718*x^3 - 0.655*x^4};
  \addlegendentry{Empírico (ec.~\ref{eq:poly})}
  % autorrotación ideal vi = -Vc
  \addplot[black!60, dashed, domain=-3.6:0, samples=2] {-x};
  \addlegendentry{Autorrotación ideal: $\vi=-V_c$}
  % punto de operación
  \addplot[only marks, mark=*, mark size=3pt, cfreno] coordinates {(-1.826,1.808)};
  \node[cfreno, font=\footnotesize, anchor=west, align=left, fill=white, inner sep=1pt] at (axis cs:-1.1,3.05)
     {Punto de diseño\\ $V/\vh=1{,}83$, $\CR\approx1{,}2$};
  \draw[cfreno, -{Latex}] (axis cs:-1.1,2.95) -- (axis cs:-1.78,1.88);
  % rótulos de zonas
  \node[font=\scriptsize, align=center, cfijo] at (axis cs:-3.0,0.3) {Molino\\ frenante};
  \node[font=\scriptsize, align=center, cfreno] at (axis cs:-1.7,0.3) {Estela\\ turbulenta};
  \node[font=\scriptsize, align=center, ccuerda!70!black] at (axis cs:-0.7,0.3) {Anillo de\\ vórtices};
  \node[font=\scriptsize, align=center] at (axis cs:0.75,0.3) {Ascenso};
\end{axis}
\end{tikzpicture}
\caption{Velocidad inducida en vuelo axial. La teoría de cantidad de movimiento vale en
ascenso y en el estado de molino frenante; entre ambos se usa la curva empírica. El rotor
libre opera cerca de la autorrotación ideal, donde el flujo neto a través del disco es casi
nulo ($V_c+\vi\approx0$).}
\label{fig:estados}
\end{figure}

\subsection{Coeficiente equivalente del rotor}
Como $T=2\rho A_R\vh^2$, el coeficiente de resistencia equivalente del rotor que baja a $V$ es
\begin{equation}
  \CR = \frac{T}{\tfrac12\rho V^2 A_R} = 4\left(\frac{\vh}{V}\right)^2 .
  \label{eq:CR}
\end{equation}
La autorrotación ideal (sin pérdidas de perfil) ocurre donde la curva empírica corta la
recta $\vi=-V_c$: con la ecuación~\eqref{eq:poly} eso sucede en $V/\vh\approx1{,}8$, es
decir $\CR\approx1{,}2$--$1{,}3$. Las pérdidas por resistencia de perfil y de punta
llevan los rotores reales a $V/\vh\approx1{,}8$--$2{,}0$, o sea
$\CR\approx1{,}0$--$1{,}2$~\cite{leishman2006,johnson1980}. Este es el
\textbf{límite físico}: ningún perfil lleva el rotor a frenar mucho más que un disco
macizo de $\CR\approx1{,}3$, por lo que la forma de mejorar la caída es el área, no el
perfil.

%=====================================================================
\section{Teoría de elemento de pala}\label{sec:bet}

\subsection{Triángulo de velocidades}

\begin{figure}[H]
\centering
\begin{tikzpicture}[scale=0.95]
  \draw[eje] (-7.6,0) -- (5.4,0) node[right, font=\footnotesize, black] {plano de giro};
  \begin{scope}[rotate=6]
    \draw[rot, line width=1pt] (-3,0) .. controls (-1.8,0.62) and (1.2,0.60) .. (3,0.04)
         .. controls (1.2,0.40) and (-1.8,0.42) .. (-3,0) -- cycle;
    \draw[cgris, densely dotted] (-4.2,0) -- (4.6,0) node[right, font=\footnotesize, black] {cuerda};
  \end{scope}
  \node[font=\footnotesize, crot, anchor=south east] at (-3.1,0.15) {borde de ataque};
  \draw[cgris] (4.3,0) arc (0:6:4.3);
  \node[font=\footnotesize, anchor=west] at (4.35,0.2) {$|\theta|$};
  % viento relativo
  \draw[flecha, cfijo] (-7.4,-1.32) -- (-3.2,-0.40);
  \node[font=\small, cfijo, anchor=south east] at (-5.2,-0.85) {$U$};
  \draw[flecha, cfijo!55] (-7.4,-1.32) -- (-3.2,-1.32) node[midway, below, font=\small] {$U_T=\Omega r$};
  \draw[flecha, cfijo!55] (-3.2,-1.32) -- (-3.2,-0.42) node[midway, right, font=\small] {$U_P=V-v$};
  \draw[cfijo] (-6.0,-1.32) arc (0:12.4:1.4);
  \node[font=\small, cfijo] at (-5.75,-1.13) {$\phi$};
  % fuerzas
  \coordinate (P) at (-1.5,0.12);
  \draw[flecha, ccuerda] (P) -- ++(102:3.4) node[above left, font=\small] {$dL$ (perpendicular a $U$)};
  \draw[flecha, cfreno] (P) -- ++(12.4:1.4) node[above right, font=\small] {$dD$};
  \draw[flecha, black!70, densely dashed] (P) -- ++(0,3.4) node[above right, font=\small] {$dT$};
  \draw[flecha, black!70, densely dashed] ($(P)+(0,1.6)$) -- ++(-1.1,0) node[left, font=\small] {$dF_Q$ impulsa};
  \fill (P) circle (1.5pt);
  \draw[-{Latex}, crot, line width=1pt] (2.6,3.2) -- (0.6,3.2) node[midway, above, font=\footnotesize] {movimiento de la pala};
  \node[font=\footnotesize, anchor=west] at (0.3,-1.3) {$\alpha=\phi+\theta$, \;$\theta<0$ (borde de ataque abajo)};
\end{tikzpicture}

\caption{Triángulo de velocidades y fuerzas en una sección de pala de un rotor libre en
descenso. El aire llega desde abajo y desde adelante; la sustentación queda inclinada
hacia adelante y es la que impulsa el giro.}
\label{fig:triangulo}
\end{figure}

En una sección a radio $r$, el aire llega con la componente tangencial $U_T=\Omega r$ y la
perpendicular $U_P=V-v$ (hacia arriba). El ángulo de flujo y el de ataque son
\begin{equation}
  \phi=\arctan\frac{U_P}{U_T}, \qquad \alpha=\phi+\theta ,
\end{equation}
con $\theta$ el paso geométrico (negativo si el borde de ataque baja). Las fuerzas de las
$B$ palas en un anillo $dr$ son
\begin{align}
  dT &= B\,\tfrac12\rho U^2 c\,\bigl(C_l\cos\phi + C_d\sin\phi\bigr)\,dr,\label{eq:dT}\\
  dQ &= B\,\tfrac12\rho U^2 c\,\bigl(C_l\sin\phi - C_d\cos\phi\bigr)\,r\,dr.\label{eq:dQ}
\end{align}

\subsection{Condición de autorrotación}
El rotor está en equilibrio cuando el par total es nulo: $Q=\int_{e}^{R}dQ=0$. En cada
sección, el término $C_l\sin\phi$ impulsa y $C_d\cos\phi$ frena; la sección es neutra si
\begin{equation}
  \tan\phi = \frac{C_d}{C_l}. \label{eq:autorrot}
\end{equation}
Un perfil con buena relación $C_l/C_d$ autorrota con $\phi$ pequeño, es decir con $U_T$
grande frente a $U_P$: gira rápido. A lo largo de la pala aparecen tres zonas
(figura~\ref{fig:zonas}): cerca de la raíz $\phi$ es grande, el perfil está en pérdida y
frena; en la zona media la sustentación inclinada impulsa; cerca de la punta $\phi$ es
pequeño y la sección es arrastrada~\cite{leishman2006}.

\begin{figure}[H]
\centering
\begin{tikzpicture}[x=0.65mm, y=1mm]
  \fill[cfreno!25] (44,0) rectangle (75,8);
  \fill[ccuerda!30] (75,0) rectangle (150,8);
  \fill[cfijo!20] (150,0) rectangle (200,8);
  \draw (0,0) rectangle (200,8);
  \fill[ffijo] (0,0) rectangle (44,8);
  \node[font=\scriptsize] at (22,4) {cubo + bisagra};
  \node[font=\scriptsize] at (59.5,4) {pérdida};
  \node[font=\scriptsize] at (112.5,4) {impulsora ($dQ>0$)};
  \node[font=\scriptsize] at (175,4) {impulsada ($dQ<0$)};
  \foreach \r/\l in {0/0,44/e,100/0{,}5R,150/0{,}75R,200/R}{
     \draw (\r,0) -- (\r,-1.5) node[below, font=\scriptsize] {$\l$};}
\end{tikzpicture}
\caption{Zonas típicas de una pala en autorrotación vertical (posiciones indicativas; la
frontera exacta depende del paso y del perfil).}
\label{fig:zonas}
\end{figure}

\subsection{Velocidad de giro de equilibrio}
Integrando~\eqref{eq:dT} con $\phi$ pequeño y $C_l$ medio constante,
\begin{equation}
  T \approx B\,\tfrac12\rho\,\Omega^2 c\,\bar C_l\,\frac{R^3-e^3}{3}
  \quad\Rightarrow\quad
  \Omega \approx \sqrt{\frac{6\,T}{B\rho\,c\,\bar C_l\,(R^3-e^3)}} .
  \label{eq:omega}
\end{equation}
El empuje $T$ lo fija el equilibrio vertical, así que \textbf{la velocidad de giro baja
con la raíz de $c\,\bar C_l$}: un perfil de alta sustentación y más cuerda gira más lento
con la misma velocidad de caída. Es el argumento del capítulo~\ref{cap:dim} para elegir
una placa curvada.

\subsection{Par de arranque}\label{sec:arranque-teoria}
Con el rotor quieto ($\Omega=0$) el aire llega perpendicular al plano de giro:
$\phi=\ang{90}$ y $\alpha=\ang{90}+\theta$. Una placa a gran ángulo de ataque genera una
fuerza normal $C_n\approx2\sin\alpha$~\cite{hoerner1965}, con sustentación
$C_l=C_n\cos\alpha$. El par de arranque es entonces
\begin{equation}
  Q_0 = B\,\tfrac12\rho V^2\,c\,C_l(\ang{90}+\theta)\,\frac{R^2-e^2}{2},
  \qquad C_l(\ang{90}+\theta)\approx -\sin 2\theta .
  \label{eq:Q0}
\end{equation}
\textbf{El signo lo decide el paso}: con $\theta<0$ (borde de ataque hacia abajo) la
fuerza normal tiene una componente hacia adelante y el rotor arranca solo; con
$\theta>0$ el par es negativo y el rotor no arranca, o arranca al revés. Esta es la razón
física para fijar el paso en $\theta\le0$.

%=====================================================================
\section{Aerodinámica a bajo número de Reynolds}\label{sec:lowre}
El número de Reynolds de la sección es $Re=U c/\nu$. Con cuerdas de
\SIrange{35}{50}{mm} y velocidades de \SIrange{10}{50}{m/s}, las palas trabajan entre
$Re\approx\num{2e4}$ y $\num{1.2e5}$.

\begin{itemize}
  \item Por debajo de $Re\approx\num{1e5}$, la capa límite laminar se separa y forma una
        \textbf{burbuja de separación laminar} que aumenta mucho la resistencia de los
        perfiles gruesos~\cite{lissaman1983,mueller2003}.
  \item Las \textbf{placas finas curvadas} no dependen de una transición controlada: el
        borde de ataque afilado fija la separación y la curvatura recupera la
        sustentación. Schmitz~\cite{schmitz1942} midió su ventaja en aeromodelos;
        Pelletier y Mueller~\cite{pelletier2000} y Okamoto y colaboradores~\cite{okamoto1996}
        confirmaron que, a $Re\sim10^4$--$10^5$, superan a los perfiles convencionales.
  \item Una curvatura de \SIrange{4}{8}{\percent} da $C_{l,\max}\approx1{,}1$--$1{,}4$ con
        relación $C_l/C_d$ de 10--20 en ese rango~\cite{pelletier2000,mueller2003,selig1995}.
\end{itemize}

\begin{figure}[H]
\centering
\begin{tikzpicture}
\begin{axis}[width=0.95\textwidth, height=8.5cm,
  xmin=-0.02, xmax=1.02, ymin=-0.19, ymax=0.62, ytick=\empty, xtick={0,0.25,0.5,0.75,1},
  xlabel={$x/c$}, grid=none, legend columns=2, legend style={font=\scriptsize, at={(0.5,-0.2)}, anchor=north,
  cells={anchor=west}}]
  % placa plana (desplazada +0.50)
  \addplot[black, thick] coordinates {(0,0.55) (1,0.55)};
  \addlegendentry{Placa plana 3\,\%}
  % placa curvada 7% (desplazada +0.30)
  \addplot[crot, very thick, domain=0:1, samples=60]
     {sqrt(1.8207^2-(x-0.5)^2)-(1.8207-0.07)+0.40};
  \addlegendentry{Placa curvada 7\,\% (elegida)}
  % fondo plano tipo Clark Y (desplazada +0.12)
  \addplot[ccuerda, thick, domain=0:1, samples=80]
     {0.22 + 0.6*(0.2969*sqrt(x)-0.1260*x-0.3516*x^2+0.2843*x^3-0.1036*x^4)};
  \addplot[ccuerda, thick, forget plot] coordinates {(0,0.22) (1,0.22)};
  \addlegendentry{Fondo plano (tipo Clark Y)}
  % NACA 6412 (desplazada -0.10)
  \addplot[cfijo, thick, domain=0:180, samples=90, variable=\t]
    ({0.5*(1-cos(t))},
     {0.03 + (0.5*(1-cos(t))<0.4 ? 0.06/0.16*(0.8*0.5*(1-cos(t))-(0.5*(1-cos(t)))^2)
                                  : 0.06/0.36*(0.2+0.8*0.5*(1-cos(t))-(0.5*(1-cos(t)))^2))
      + 5*0.12*(0.2969*sqrt(0.5*(1-cos(t)))-0.1260*0.5*(1-cos(t))-0.3516*(0.5*(1-cos(t)))^2
      +0.2843*(0.5*(1-cos(t)))^3-0.1015*(0.5*(1-cos(t)))^4)});
  \addplot[cfijo, thick, domain=0:180, samples=90, variable=\t, forget plot]
    ({0.5*(1-cos(t))},
     {0.03 + (0.5*(1-cos(t))<0.4 ? 0.06/0.16*(0.8*0.5*(1-cos(t))-(0.5*(1-cos(t)))^2)
                                  : 0.06/0.36*(0.2+0.8*0.5*(1-cos(t))-(0.5*(1-cos(t)))^2))
      - 5*0.12*(0.2969*sqrt(0.5*(1-cos(t)))-0.1260*0.5*(1-cos(t))-0.3516*(0.5*(1-cos(t)))^2
      +0.2843*(0.5*(1-cos(t)))^3-0.1015*(0.5*(1-cos(t)))^4)});
  \addlegendentry{NACA 6412 (Ventura 2022)}
  % NACA 0012 (desplazada -0.27)
  \addplot[cgris, thick, domain=0:180, samples=90, variable=\t]
    ({0.5*(1-cos(t))}, {-0.12 + 0.6*(0.2969*sqrt(0.5*(1-cos(t)))-0.1260*0.5*(1-cos(t))
      -0.3516*(0.5*(1-cos(t)))^2+0.2843*(0.5*(1-cos(t)))^3-0.1015*(0.5*(1-cos(t)))^4)});
  \addplot[cgris, thick, domain=0:180, samples=90, variable=\t, forget plot]
    ({0.5*(1-cos(t))}, {-0.12 - 0.6*(0.2969*sqrt(0.5*(1-cos(t)))-0.1260*0.5*(1-cos(t))
      -0.3516*(0.5*(1-cos(t)))^2+0.2843*(0.5*(1-cos(t)))^3-0.1015*(0.5*(1-cos(t)))^4)});
  \addlegendentry{NACA 0012 (helicópteros)}
\end{axis}
\end{tikzpicture}
\caption{Secciones comparadas, desplazadas en vertical para verlas por separado. La
escala vertical está ampliada unas 2,5 veces respecto de la horizontal. El espesor de la
placa curvada se dibuja como una línea porque es el 1,1\,\% de la cuerda.}
\label{fig:perfiles}
\end{figure}

\begin{table}[H]
\centering\small
\caption{Características aproximadas a $Re\approx\num{5e4}$--$\num{1e5}$ (orden de
magnitud de la literatura~\cite{pelletier2000,mueller2003,selig1995,okamoto1996}; deben
verificarse en ensayo).}\label{tab:polares}
\begin{tabular}{@{}lcccl@{}}
\toprule
\textbf{Sección} & $\boldsymbol{C_{l,\max}}$ & $\boldsymbol{(C_l/C_d)_{\max}}$ & $\boldsymbol{C_l}$ \textbf{de trabajo} & \textbf{Fabricación}\\\midrule
Placa plana 3\,\%            & 0,8--0,9 & 6--10  & $\approx0{,}3$ & Muy fácil\\
Fondo plano (Clark Y aprox.) & 1,0--1,2 & 10--15 & $\approx0{,}6$ & Lija y plantilla\\
Placa curvada 4\,\%          & 1,0--1,2 & 12--18 & $\approx0{,}7$ & Molde\\
\textbf{Placa curvada 7\,\%} & \textbf{1,2--1,4} & \textbf{12--20} & $\boldsymbol{\approx0{,}9}$ & \textbf{Molde}\\
NACA 6412                    & 1,1--1,3 & 10--20 & $\approx0{,}8$ & Impresión o fresado\\
NACA 0012                    & 0,7--0,9 & 5--10  & $\approx0{,}4$ & Impresión o fresado\\
\bottomrule
\end{tabular}
\end{table}

%=====================================================================
\section{Dinámica de batimiento}\label{sec:batimiento}

\begin{figure}[H]
\centering
\begin{tikzpicture}[scale=1.0]
  \draw[fijo] (-0.6,-0.8) rectangle (0,0.6);
  \node[font=\scriptsize, cfijo, rotate=90] at (-0.3,-0.1) {cubo};
  \draw[eje] (-0.3,-1.2) -- (-0.3,2.4) node[above, font=\footnotesize, black] {eje de giro};
  \fill[black] (0.25,0) circle (2pt); \node[font=\footnotesize, anchor=north west] at (0.35,-0.25) {bisagra};
  \draw[eje] (0.25,0) -- (8.2,0);
  \begin{scope}[rotate around={7:(0.25,0)}]
    \draw[rot] (0.25,-0.07) rectangle (8,0.07);
    \draw[flecha, ccuerda] (5.6,0.1) -- ++(0,1.5) node[above, font=\small] {$L$ (sustentación)};
    \draw[flecha, cfreno] (4.1,0) -- ++(1.6,0);
    \fill[black] (4.1,0) circle (1.6pt);
    \draw[flecha, cgris] (4.1,-0.05) -- ++(0,-0.9) node[below, font=\small] {$m_b g$};
  \end{scope}
  \node[cfreno, font=\small, anchor=north] at (6.4,0.35) {$F_c = m_b r_g \Omega^2$};
  \draw[cgris] (2.4,0) arc (0:7:2.15);
  \node[font=\small] at (2.75,0.17) {$\beta_0$};
  \draw[cota] (0.25,-1.6) -- (8.0,-1.6) node[midway, below, font=\footnotesize] {$L_b = R-e$};
  \draw[cota] (-0.3,-2.2) -- (0.25,-2.2) node[midway, below, font=\footnotesize] {$e$};
\end{tikzpicture}
\caption{Equilibrio de una pala articulada: la sustentación la levanta y la fuerza
centrífuga la lleva hacia el plano de giro. El resultado es un ángulo de conicidad
$\beta_0$ de pocos grados.}
\label{fig:batimiento}
\end{figure}

Cada pala se une al cubo con una bisagra de batimiento a radio $e$, como en el rotor de
Cierva. Para ángulos pequeños, el equilibrio de momentos en la bisagra da
\begin{equation}
  \beta_0 = \frac{M_L - M_W}{m_b\,\Omega^2 \,\overline{r(r-e)}},\qquad
  \overline{r(r-e)} = e\frac{L_b}{2}+\frac{L_b^2}{3},
  \label{eq:conicidad}
\end{equation}
con $M_L$ el momento de la sustentación de la pala respecto de la bisagra, $M_W$ el de su
peso y $m_b$ su masa (pala uniforme). Como la sustentación y la fuerza centrífuga crecen
ambas con $\Omega^2$, la conicidad casi no depende de la velocidad de giro, salvo en el
arranque.

La relación entre fuerzas aerodinámicas e inerciales la mide el \textbf{número de Lock}
\begin{equation}
  \gamma = \frac{\rho\, a\, c\, R^4}{I_b}, \qquad I_b \approx \frac{m_b L_b^2}{3},
  \label{eq:lock}
\end{equation}
con $a$ la pendiente de sustentación. En helicópteros $\gamma\approx5$--10; valores
mayores indican palas muy livianas, con mucha amortiguación aerodinámica pero conicidad
grande~\cite{johnson1980,leishman2006}. La frecuencia natural de batimiento de una pala
articulada con excentricidad $e$ es
\begin{equation}
  \frac{\nu_\beta}{\Omega} = \sqrt{1+\frac{3}{2}\frac{e}{L_b}} .
  \label{eq:nubeta}
\end{equation}

%=====================================================================
\section{Paracaídas de frenado (drogue)}\label{sec:teoria-drogue}
El drogue se dimensiona con la ecuación~\eqref{eq:equilibrio}. Knacke~\cite{knacke1992}
da coeficientes, referidos al área nominal, de \numrange{0.75}{0.80} para paracaídas
circulares planos y valores similares, algo menores, para los de cruz. La carga de apertura
se estima con
\begin{equation}
  F_o = C_x\, X_1\, q\, C_D S, \qquad q=\tfrac12\rho V_{\text{ap}}^2 ,
  \label{eq:apertura}
\end{equation}
donde $C_x\approx1{,}5$--1,8 es el factor de choque a masa infinita y $X_1\le1$ el factor
de reducción por masa finita~\cite{knacke1992}. Se toma $C_x=1{,}5$ y $X_1=1$
(conservador).

\section{Resistencia del cuerpo}
Un cilindro romo moviéndose según su eje tiene $C_D\approx0{,}8$--1,0 para relaciones de
aspecto de 2 a 3~\cite{hoerner1965}. Se usa $C_{D,c}=0{,}8$ sobre el área frontal de
Ø\,\SI{95}{mm}. Su aporte es menor al 3\,\% de la resistencia total en la segunda etapa.
